\chapter{The Piston-Cylinder: the Heat-Engine Problem}\label{ch:piston}

\small
\begin{quote}
The steam engine having furnished us with a means for converting heat
into motive power\ldots\ the idea of establishing theoretically some
fixed relation between a quantity of heat and the quantity of work
which it can possibly produce, from which relation conclusions
regarding the nature of heat itself may be deduced, naturally
presents itself.
\normalfont\hfill --- R.~Clausius, 1851.
\end{quote}
\normalsize

\section{The piston-cylinder as the heat-engine analogue of Joule}

The Joule experiment of Chapter~\ref{ch:joule} ran a pressure differential
between two fixed-wall chambers to mechanical equilibrium and observed
the residual temperature/density gap: a one-shot dissipation pattern with
no net work extracted. The companion foundational problem of
thermodynamics replaces the right-hand fixed wall of the Joule chamber
with a \emph{movable piston} and connects the piston to an external
mechanical load. Gas now does work on the load as it expands, and the
question becomes how much of the internal-energy decrease in the
gas reservoir reaches the load as work versus how much is irreversibly
dissipated as heat in the receiving region.

We treat two cases. The \emph{smooth} case is a single chamber filled
with hot dense gas, closed on three sides and with a piston as the
fourth wall. The gas expands quasi-statically against the load; in the
ideal limit of slow piston motion this approaches reversible expansion
with $\eta = W/\Delta U \to 1$. The \emph{Joule-coupled} case inserts a
constricted wall (a thin barrier with a small channel opening) between
the gas reservoir and the piston. The gas now jets through the
constriction, decelerates against the gas already in the receiving
region, and the recompression-and-turbulent-dissipation pattern of
Joule's experiment plays out --- but now with a piston on the far side
that captures whatever pressure survives. The efficiency drops below
the smooth limit by the amount of kinetic energy that turbulent
recompression converts back to heat rather than to useful pressure on
the piston. The two cases together display the fundamental engine
trade-off in the cleanest possible computational form.

\section{The smooth piston-cylinder}\label{sec:piston-smooth}

\paragraph{Configuration.}
A square domain $[0, 1]^2$ holds gas of initial density $\rho = 10$ and
internal energy $e = 10$ (so $p = \gamma e = 1.0$ with $\gamma = 0.1$)
to the left of a movable piston initially positioned at $x = 0.5$. To
the right of the piston lies an external region held at a fixed load
pressure $P_{\rm ext} = 0.5$. The top and bottom walls of the chamber
are zero-flux/free-slip, the left wall is closed, the right boundary
beyond the piston is the load.

\paragraph{Dynamics.}
Newton's second law for the piston, with $M$ the piston mass per unit
$y$-length and a friction coefficient $c$ for piston-rail damping, reads
\begin{equation}\label{eq:piston-newton}
M\,\frac{\rd v_p}{\rd t} = \int_{0}^{1}\!\bigl[\,p_{\rm gas}(x_p^-, y) - P_{\rm ext}\,\bigr]\,\rd y \;-\; c\,v_p,
\quad \frac{\rd x_p}{\rd t} = v_p,
\end{equation}
where $p_{\rm gas}(x_p^-, y) = \gamma\,e(x_p^-, y)$ is the local
gas pressure on the piston's gas-side face. The gas occupies the
domain $x \in [0, x_p)$ and evolves by the compressible
Navier--Stokes equations of Chapter~\ref{ch:joule}, with a no-slip
ghost at the piston face. Work delivered to the load per unit time is
\begin{equation}\label{eq:piston-power}
\dot W = P_{\rm ext} \cdot v_p \cdot \mathrm{Area},
\end{equation}
and the heat-engine efficiency is
\begin{equation}\label{eq:piston-eta}
\eta = \frac{W}{\Delta U} = \frac{\int_0^t \dot W\,\rd s}{U(0) - U(t)}.
\end{equation}

\paragraph{Behaviour and expected result.}
The smooth piston-cylinder is the cleanest computational instance of
quasi-static isobaric expansion. With slow piston motion (small $v_p$
compared with the acoustic speed $c_s = \sqrt{\gamma e/\rho}$ in the
gas), the gas expands without significant kinetic-energy generation;
the work-delivered counter and the heat-lost counter grow together;
and $\eta$ approaches the reversible limit. Faster piston motion
generates acoustic waves in the gas region that bounce between the
left wall and the piston face, with each acoustic round-trip a small
turbulent-dissipation loss; $\eta$ falls below 1 by an amount
proportional to the piston Mach number. The smooth case therefore
isolates the fundamental work-extraction mechanism, with irreversibility
present only as a small correction proportional to the piston speed.

The accompanying simulator \texttt{heat\_engine.html} (Appendix~A
Gallery) implements~\eqref{eq:piston-newton}--\eqref{eq:piston-eta} for
$N = 200$, $\gamma = 0.1$, $P_{\rm ext} = 0.5$, $M = 0.05$,
$c = 1.5$. The piston moves from $x_p = 0.5$ outward over $\sim 10\,$s
of wall-clock time, with $\eta$ reaching $\sim 0.7$--$0.9$ depending on
the piston-Mach-number regime.

\section{The Joule-coupled piston-cylinder}\label{sec:piston-joule}

\paragraph{Configuration.}
The same square domain holds a high-pressure gas reservoir on the left
($i < I_w$ in grid units, $\rho = 10$, $e = 10$), a thin barrier wall at
$i \in [I_w, I_w + w_{\rm wall}]$ with a small channel opening at
$j \in [J_0, J_1]$, a receiving region between the barrier and the
piston ($I_w + w_{\rm wall} < i < x_p$, initially $\rho = 1$, $e = 1$),
and the movable piston at $i = x_p$. The external load is held at
$P_{\rm ext} = 0.1$, matching the initial pressure in the receiving
region so the piston starts in mechanical equilibrium and is pushed
outward only after the channel inflow has compressed the receiving
gas.

\paragraph{Dynamics.}
The piston obeys the same Newton's-law update~\eqref{eq:piston-newton}
with $p_{\rm gas}$ now the local pressure in the receiving region
just to the left of the piston face. The gas equations are unchanged;
the barrier is a frozen no-slip wall except in the channel opening,
where gas can flow freely. The barrier presents the classical
Joule constriction to the gas: the high reservoir pressure drives a
jet through the channel, the jet decelerates in the receiving region,
and the kinetic-energy-to-heat conversion of turbulent
recompression dissipates a fraction of the released internal energy
into heat in the receiving region rather than into useful pressure
against the piston.

\paragraph{Two energy accounts.}
The 1st Law of thermodynamics for this system says
\begin{equation}\label{eq:piston-1st-law}
\underbrace{\Delta U_{\rm reservoir}}_{\text{heat lost from hot side}}
\;=\;
\underbrace{W}_{\text{useful work}}
\;+\;
\underbrace{\Delta U_{\rm receiving}}_{\text{heat dumped in cold side}}
\;+\;
\underbrace{K_{\rm channel}}_{\text{kinetic energy in flight}}.
\end{equation}
The transient $K_{\rm channel}$ goes to zero at long times as the gas
equilibrates; in the steady state, all of the released $\Delta U_{\rm reservoir}$
ends up either as work $W$ or as $\Delta U_{\rm receiving}$ (waste heat).
The 2nd Law tells us which split is realised: the size of
$\Delta U_{\rm receiving}$ is set by the dynamic
turbulent-dissipation pattern in the recompression process, exactly
as in the Joule experiment of Chapter~\ref{ch:joule}. The engine
efficiency follows as
\begin{equation}\label{eq:piston-eta-joule}
\eta \;=\; \frac{W}{\Delta U_{\rm reservoir}} \;=\; 1 - \frac{\Delta U_{\rm receiving}}{\Delta U_{\rm reservoir}}\;-\;\frac{K_{\rm channel}}{\Delta U_{\rm reservoir}},
\end{equation}
and in the steady state $\eta < 1$ by exactly the fraction of energy
that the Joule turbulent-recompression channel dissipates as heat.

\paragraph{One-shot versus sustained operation.}
With the left wall closed (finite reservoir), the simulation runs as a
\emph{one-shot} cycle: the reservoir depletes, the piston receives a
brief impulsive load from the channel-driven shock and a slower
follow-on pressure as the gas equilibrates, and the engine extracts a
single power stroke of work. This is the elementary unit of a
reciprocating closed-cycle heat engine (Otto, Diesel, Stirling).
With the left wall replaced by a Dirichlet inflow boundary maintaining
the reservoir state at $\rho = 10$, $e = 10$ indefinitely, the
simulation runs as a \emph{sustained} cycle: gas keeps jetting through
the channel, the piston reaches a quasi-steady velocity, and power is
delivered continuously. This is the elementary unit of an open-cycle
flow engine (gas turbine, jet engine). The two operating modes are
selected in the simulator \texttt{heat\_engine\_with\_channel.html}
(Appendix~A) by the choice of left-wall boundary condition: zero-gradient
mirror for one-shot, Dirichlet inflow for sustained.

\section{Joule and the piston as duals}\label{sec:joule-piston-duals}

The two central problems are duals of each other:
\begin{itemize}
\item \textbf{Joule's experiment} (Chapter~\ref{ch:joule}) lets the
expansion-and-recompression dynamics run to mechanical equilibrium
with no work extracted, leaving a residual temperature/density gap
between the chambers that is the dynamic-dissipation signature. With
external work added cyclically by a compressor, this becomes a
\textbf{refrigerator}: the residual asymmetry is held against
diffusive relaxation by continuously re-energising the gas, and the
asymmetric temperature gradient is exploited to draw heat from cold
side and dump it to hot side.

\item \textbf{The piston-cylinder} (this chapter) lets the same
expansion-and-recompression dynamics run with a movable wall as the
right boundary, extracting useful work to an external load. With
external heat added cyclically by a hot reservoir at the inflow, this
becomes a \textbf{heat engine}: the released internal energy is split
between work delivered to the load and waste heat dumped to a cold
side, with the efficiency $\eta$ set by the same turbulent-dissipation
fraction.
\end{itemize}
Both problems are governed by the same compressible
Navier--Stokes equations, the same RNS finite-element discretisation,
and the same 2nd-Law-as-finite-precision reading of irreversibility.
The framework of this book treats them as the two foundational
single-shot configurations of which all cyclic engines, pumps,
refrigerators, and heat-recovery devices of Vol VI Applications are
built; everything else in thermodynamics is either a cyclic repetition
of these or a multi-stage combination of them.

\section{What this chapter establishes}

For the rest of the book the reader should hold three observations
from this chapter and the previous one together:
\begin{enumerate}
\item Pressure-driven flow through a constriction converts internal
energy into kinetic energy; recompression on the receiving side
converts a fraction of that kinetic energy back into heat
irreversibly; the size of that fraction is the dynamic content of the
2nd Law and is what classical equilibrium thermodynamics cannot
predict.

\item A piston in the receiving region intercepts a fraction of the
arriving pressure as useful work; the larger that fraction, the
smaller the recompression dissipation, the higher the engine
efficiency. The efficiency is bounded above by the reversible smooth-
expansion limit, and the irreversibility cost is the dynamical content
of the heat-engine inequality $\eta \le 1$.

\item Both the refrigerator and the heat engine are cyclic devices
built from a single foundational dynamic process: the
expansion-through-a-channel-followed-by-turbulent-recompression cycle.
Refrigerators run it as the heat-pumping cycle with work input;
heat engines run it as the work-output cycle with heat input. The
same simulation infrastructure produces both.
\end{enumerate}

